\documentclass{ximera}[font=12pt]  % I don't think font sizing works here

\title{Fractions: GCF}   % Use xmlatex all -s to compile when SERVE refuses to work.
\author{Kelly Stady}
\license{CC: 0}         % replace with an appropriate license, or set it in xmPreamble

\begin{document}


\begin{abstract}

\end{abstract}

\maketitle

\label{xim:FracsGCFSimplify}


\section*{Greatest Common Factor (GCF)}

\begin{definition}
The \textbf{greatest common factor} is the largest whole number that divides evenly into two or more numbers.  Equivalently, it is
the largest whole number that is a factor of each number.

\end{definition}


\begin{problem} \textcolor{glaucous!55!black}{\textbf{True or False}}:  The GCF will always be less than or equal to the smallest number in the group.

\begin{multipleChoice}
\choice[correct]{True.}
\choice{False.}
\end{multipleChoice}

\begin{feedback}[correct]
That's right!  Since the GCF evenly divides into each number, it can't be larger than the smallest number.
\end{feedback}

\end{problem}



\begin{procedure}[\textbf{Finding the GCF Using Lists of Factors}]

List the factors of each number and look for the largest number that appears in each list.  \textbf{Alternatively}, 
only list the factors of the smallest number then look for the largest factor that evenly divides into the other number(s).

\end{procedure}


\begin{example}

To find the GCF of 12 and 36, we begin by listing the factors of each number.


\qquad \textbf{Factors of 12:} 1, 2, 3, 4, 6, 12

\qquad \textbf{Factors of 36:} 1, 3, 4, 6, 9, 12, 18, 36


The \textit{common factors} of 12 and 36 are 1, 3, 4, 6, and 12.  So, the \textbf{greatest common factor} is 12.  
Notice that it is more efficient to list the factors of 12 and look for the largest number that evenly divides into 36.  

\end{example}


INSERT VIDEO SHOWING Alternative IDEA


\begin{problem}
Find the GCF of 15 and 20 by listing the factors of both numbers or just the number 15.  Use whichever technique you prefer.

\[ \textbf{GCF} = \answer{5} \]

\end{problem}


\begin{problem}
Find the GCF of 24 and 60 by listing the factors of both numbers or just the number 24.  Use whichever technique you prefer.

\[ \textbf{GCF} = \answer{12} \]

\end{problem}


\begin{procedure}[\textbf{Finding the GCF Using Prime Factorizations}]
    
    Find the prime factorization of the numbers and multiply the common prime factors.  
    
\end{procedure}

\begin{example}

 To find the GCF of 12 and 36, using this approach, we begin with the prime factorization of both numbers.

\begin{eqnarray*}
12 &=& 2 \cdot 2 \cdot 2 \cdot 3 \\
   & & \\
36 &=& 2 \cdot 2 \cdot 3 \cdot 3 \\
\end{eqnarray*}

The common factors are 2 and 3.  There are two factors of 2 and one factor of 3 in \textit{both} factorizations, 
so the GCF = $2 \cdot 2 \cdot 3$ which equals 12.
\end{example}

\begin{problem}
Find the GCF of 60 and 75 using their prime factorizations shown below.

\begin{eqnarray*}
60 &=& 2 \cdot 2 \cdot 3 \cdot 5 \\
   & & \\
75 &=& 3 \cdot 5 \cdot 3 \\
\end{eqnarray*}

\textbf{GCF} = $\answer{15}$

\begin{feedback}
Excellent!  The common factors are 3 and 5.  There is one factor of 3 and one factor of 5 in \textbf{both} factorizations, so the GCF = $3 \cdot 5$ which equals 15.
\end{feedback}

\end{problem}


\begin{problem}
Find the GCF of 85 and 150 using their prime factorizations shown below.


\begin{eqnarray*}
85 &=& 5 \cdot 17 \\
   & & \\
150 &=& 2 \cdot 3 \cdot 5 \cdot 5 \\
\end{eqnarray*}


\textbf{GCF} = $\answer{5}$

\end{problem}


INSERT VIDEO:  Finding the GCF of 3 numbers!!!!!

The GCF can be used to solve problems involving groups of things that need to be divided efficiently, with no waste.  Use what you've learned
about finding the GCF to solve the following application problems.

\begin{problem}
  You are creating identical goodie bags for a school event. You have \textbf{18} Airheads, \textbf{24} Tootsie Rolls and
  \textbf{36} Kit-Kats.  What is the \textbf{greatest} number of identical goodie bags you can make if you want to use all of the candy?

  \[ \textrm{Greatest Number of Goodie Bags} = \answer{6} \] 

  \begin{hint}
    To solve this, you need to find the Greatest Common Factor (GCF) of 18, 24, and 36. List the factors for each number and find the largest one that appears in all three lists.
  \end{hint}

  Now that you know you can make identical goodie bags with no waste, how many of each candy will be in \textbf{each} bag?

  \begin{itemize}
      \item Airheads: $\answer{3}$
      \item Tootsie Rolls: $\answer{4}$
      \item Kit-Kats: $\answer{6}$
  \end{itemize}
  
  \begin{hint}
    Divide the total number of each candy by the number of bags (6). For example, $18 \div 6 = 3$ Airheads.
  \end{hint}


\end{problem}


\begin{problem}
  A carpenter has three wooden boards of different lengths: \textbf{8} feet, \textbf{12} feet and
      \textbf{20} feet.  The carpenter wants to cut all three boards into smaller pieces of equal length,
       so that there is no scrap wood left over. What is the longest possible length for these pieces?

  \[ \textrm{Longest Length } = \answer{4} \ \textrm{feet} \] 

  \begin{hint}
    Find the Greatest Common Factor (GCF) of 8, 12, and 20. This will be the largest number that divides evenly into all three lengths.
  \end{hint}

How many pieces will the carpenter have after cutting all three boards?

  \[ \textrm{Total Pieces } = \ \answer{10} \] 

  \begin{hint}
    Divide each board's length by 4 to see how many pieces come from each:
    \begin{itemize}
        \item $8 \div 4 = 2$ pieces
        \item $12 \div 4 = 3$ pieces
        \item $20 \div 4 = 5$ pieces
    \end{itemize}
    Now, add them.
  \end{hint}
\end{problem}


\section*{Simplifying Fractions}

A fraction is said to be in \textbf{simplest form} or \textbf{lowest terms} when the
numerator and denominator have no common factors other than 1.  For example, the fraction
$\textstyle \frac{2}{5}$ is in simplest form because 2 and 5 have no common factors.  However, the fraction
$\textstyle \frac{4}{6}$ is not in simplest form because both 4 and 6 have a common factor of 2.  The
process of writing a fraction in simplest form is called \textbf{simplifying} the fraction.  We will consider two methods 
to simplify fractions.  The first method uses the GCF and the second method uses prime factorizations.

\begin{procedure}[\textbf{Simplifying Fractions Using the GCF}]

    Divide the numerator and denominator by the GCF.

\end{procedure}

\begin{example}

The fraction $\textstyle \frac{4}{12}$ can be simplified by dividing the numerator and denominator by the GCF = 4.

\begin{align*}
\frac{4}{12} &= \frac{4 \div 4}{12 \div 4} \\[1em]
             &= \frac{1}{3}
\end{align*}


We can also simplify fractions by dividing the numerator and denominator by \emph{any} common factor, however using the GCF involves just one step.
If we didn't notice that the GCF was 4, but knew that both numbers were evenly divisible by 2, we would need two steps to simplify 
the fraction $\textstyle \frac{4}{12}.$

\begin{align*}
\frac{4}{12} &= \frac{4 \div 2}{12 \div 2} \\[1em]
             &= \frac{2}{6}             \\[1em]
             &= \frac{2 \div 2}{6 \div 2} \\[1em]
             &= \frac{1}{3}
\end{align*}

\end{example}


\begin{problem}
  Simplify the fraction below using the greatest common factor, which is 12.

  \[
  \frac{24}{36} = \frac{\answer{2}}{\answer{3}}
  \]

\end{problem}



\begin{problem}
  Simplify the fraction below.

  \[ \frac{18}{45} \]

  First, find the Greatest Common Factor (GCF) of 18 and 45.

  GCF = $\answer{9}$

  Now, divide both the numerator and the denominator by the GCF to simplify:

  \[ \frac{18}{45} = \frac{\answer{2}}{\answer{5}} \]

  \begin{hint}
    List the factors for each number to find the largest one they share (GCF):
    \begin{itemize}
        \item Factors of 18: 1, 2, 3, 6, 9, 18
        \item Factors of 45: 1, 3, 5, 9, 15, 45
    \end{itemize}
  \end{hint}
\end{problem}

\begin{procedure}[\textbf{Simplyfing Fractions Using Prime Factorizations}]

Write the prime factorization of the numbers in the numerator and denominator, then cancel the common prime factors.

\end{procedure}

\begin{example}

Simplify the fraction $\textstyle \frac{20}{25}$ using prime factorizations.

\begin{align*}
\frac{20}{25} &= \frac{2 \cdot 2 \cdot 5}{5 \cdot 5} \\[1em]
              &= \frac{2 \cdot 2 \cdot \cancel{5}}{5 \cdot \cancel{5}} \\[1em]
              &= \frac{4}{5}
\end{align*}

\end{example}


INSERT VIDEO of simplifying a fraction using one of these facts.

\begin{problem}
Simplify the fraction using the prime factorizations shown below.

\begin{eqnarray*}
    42 &=& 2 \cdot 3 \cdot 7 \\
    70 &=& 2 \cdot 5 \cdot 7
\end{eqnarray*}

$\displaystyle \frac{42}{70} = \frac{\answer{3}}{\answer{5}} $

  \begin{hint}
    Write the prime factorizations in the fraction, then cancel the common prime numbers.
    \[ \frac{42}{70} = \frac{2 \cdot 3 \cdot 7}{2 \cdot 5 \cdot 7} \]
  \end{hint}

\end{problem}


\begin{problem}
Simplify the fraction using whichever method you prefer.

\[ \frac{33}{60} = \frac{\answer{11}}{\answer{20}}\]

\begin{hint}
   If you want to use prime factorizations, write the prime factorizations in the fraction, then cancel the common prime numbers.

    \[ \frac{33}{60} = \frac{3 \cdot 11}{2 \cdot 2 \cdot 3 \cdot 5} \]

  \end{hint}

  \begin{hint}
    If you want to use the Greatest Common Factor (GCF), the largest number that divides both 33 and 60 is 3.
  \end{hint}

\end{problem}



\begin{problem}
  Simplify the fraction using whichever method you prefer.


  \[ \frac{60}{126} = \frac{\answer{10}}{\answer{21}} \]

  \begin{hint}
   If you want to use prime factorizations, write the prime factorizations in the fraction, then cancel the common prime numbers.

    \[ \frac{60}{126} = \frac{2 \cdot 2 \cdot 3 \cdot 5}{2 \cdot 3 \cdot 3 \cdot 7} \]

  \end{hint}

  \begin{hint}
    If you want to use the Greatest Common Factor (GCF), the largest number that divides both 60 and 126 is 6.
  \end{hint}

\end{problem}


\begin{problem}
  A recycling center collected $54$ tons of plastic last month.  However, due to contamination (like food waste or non-recyclable items), 
  only $21$ tons could be recycled.  What \textbf{fraction} of the collected plastic was recycled?  \textbf{Write your answer in simplest form}.

  \[ \text{Fraction Recycled} = \frac{\answer{7}}{\answer{18}} \]

  \begin{hint}
    Set up the fraction as the amount recycled over the total amount collected, then simplify.
    \[ \frac{21}{54} \]
  \end{hint}

\end{problem}


\begin{problem}
An animal rescue organization saved $120$ dogs from an overcrowded shelter.  Of these dogs, $45$ were seniors who required specialized 
medical care and longer stays before they were ready for adoption.  What \textbf{fraction} of the rescued dogs were seniors?  
\textbf{Write your answer in simplest form}.

\[ \text{Fraction of Senior Dogs} = \frac{\answer{3}}{\answer{8}} \]

\begin{hint}
Write the fraction as the number of senior dogs over the total number of dogs rescued, then simplify.

\[ \frac{45}{120} \]
\end{hint}
\end{problem}

\end{document}